\chapter{Time Integration and Total Solution Cost}\label{v0ch:25}
\begin{chaptersummary}
Spatial throughput gives the cost of one call; time integration determines call counts and implicit-solve work. A common framework connects RK, staggered integration, ADER, CERK, LTS, and IMEX to stability, phase accuracy, and structure preservation. At a specified error and final time, step size, stages, solves, and communication jointly determine total cost.
\end{chaptersummary}
\learninggoals{Relate semidiscrete spectra to stability regions; distinguish symplecticity from energy conservation; analyze high-order, local-step, and stiff-splitting costs.}
\paragraph{Core Question}Why can a simulation remain slow when one spatial evaluation is fast?

\section{Method of Lines and Stability Regions}
Spatial discretization gives $M\dot{\bm u}=\bm R(\bm u)$. Each explicit RK stage applies $M^{-1}\bm R$, cheaply for local or diagonal mass matrices. For $\dot y=\lambda y$, one step is $y^{n+1}=R(z)y^n$ with $z=\Delta t\lambda$; absolute stability requires $|R(z)|\le1$. Classical RK4 has
$R(z)=1+z+z^2/2+z^3/6+z^4/24$. Spectra, nonnormality, and energy estimates jointly constrain usable steps.
For purely imaginary $\lambda$, forward Euler has $|1+\mathrm i\omega\Delta t|>1$, producing modal growth. The time integrator's stability region is therefore part of the complete energy analysis.
\section{Staggering, Symplecticity, and Phase}
Yee/leapfrog staggers suitable split systems and is conditionally stable. St\"ormer--Verlet is second-order symplectic for separable Hamiltonian systems. Implicit midpoint
\[
 \frac{\bm u^{n+1}-\bm u^n}{\Delta t}
 =A\frac{\bm u^{n+1}+\bm u^n}{2}
\]
exactly preserves a discrete quadratic energy when $A^TQ+QA=0$ with fixed SPD $Q$. For general Hamiltonian systems it is symplectic, yielding bounded long-time energy oscillations rather than exact conservation of the original Hamiltonian.
For $\lambda=\mathrm i\omega$, midpoint gives the numerical frequency
$\omega_h=2\arctan(\omega\Delta t/2)/\Delta t$, showing that amplitude stability can still accompany accumulated phase error; inexact linear solves likewise limit conservation accuracy.
\section{ADER, CERK, and LTS}
ADER constructs high-order temporal representations via element-local space--time prediction followed by interface space--time flux correction; while prediction can improve local data reuse, its net cost, stability, and order depend on the specific construction.
CERK denotes continuous explicit Runge--Kutta, providing a within-step approximation
\[
 \bm u(t_n+\theta\Delta t)
 \approx\bm u_n+\Delta t\sum_i b_i(\theta)\bm k_i,\quad0\le\theta\le1.
\]
Within-step accuracy is governed by the continuous-extension order conditions, whereas endpoint accuracy follows the standard RK order conditions, so coupling and dense output verify both orders.
LTS groups elements by size or local wave speed, maintaining conservation through a single shared time-integrated interface flux that reconciles fine-step accumulation with coarse-step correction, while interpolation, flux registers, scheduling, and load imbalance introduce overhead.
\section{IMEX and Stiff Sources}
Splitting $\dot{\bm u}=F(\bm u)+G(\bm u)$ treats propagation explicitly and stiff sources implicitly. Additive RK stages take the form
\[
 U_i=\bm u^n+\Delta t\sum_{j<i}\widetilde a_{ij}F(U_j)
                  +\Delta t\sum_{j\le i}a_{ij}G(U_j),
\]
with separate explicit and implicit endpoint weights; overall order and stability depend on the paired tableaux and coupling conditions.
In conductive Maxwell equations, taking $G_E=-\sigma E/\varepsilon$ yields the first-order split update
\[
 E^{n+1}=\frac{E^n+\Delta t\,(\curl H)^n/\varepsilon}
 {1+\Delta t\,\sigma/\varepsilon}.
\]
This illustrates inexpensive local implicit damping. Because forward Euler is unstable for pure wave propagation ($\sigma=0$), practical wave simulations pair the implicit step with an explicit tableau whose stability region covers the imaginary axis; global diffusion or coupled implicit sources, in turn, reintroduce large linear solves.
\conceptbridge{Why do high-order or implicit methods not automatically reduce total cost?}{Higher order alters stage counts and error constants, implicit methods trade stability restrictions for linear or nonlinear solves, and local time stepping adds scheduling and communication overhead. Comparing schemes at a fixed error target requires evaluating usable step sizes, total step counts, and per-step execution graphs together.}{Stability and accuracy $\longrightarrow$ usable time step $\longrightarrow$ stages, solves, and communication.}
\section{Minimal Comparisons}
Periodic sine transport compares RK4, second-order methods, and staggered wave updates, while $\dot y=\mathrm i\omega y-\sigma y$ compares explicit and IMEX stability and phase accuracy, evaluated by error at a fixed physical final time rather than $|R|$ alone.
\costnote{$T_{\rm total}=T_{\rm setup}+N_tT_{\rm step}+T_{\rm output}$, where $T_{\rm step}$ encompasses all stages, implicit solves, and communication. ADER and LTS trade local arithmetic for fewer global updates, CERK supplies time-aligned states, and interface fluxes together with the combined stability region govern conservation and stability.}
\aiexercise{Derive midpoint's phase frequency from its amplification factor. Quantify phase error after one hundred periods and compare RK4.}{Write shared flux integrals and fine-step accumulation for adjacent unequal-step elements. Explain why interpolated boundary values alone do not ensure conservation.}
\studyroute{ADER predictors, CERK order conditions, and high-order IMEX tableaux are omitted. See Volume II's time-integration chapter and \cite{hesthaven,ader,cerk,imex}.}{method of lines; symplectic integrator; ADER; continuous explicit Runge Kutta; conservative LTS; IMEX additive RK}
